\chapter{Conclusion} \label{chapter: conclusion}

In this final chapter, the focus will be to briefly review the results, and present potential further aspects 
of scientific research in that area. For this, we want to lastly distinguish
this portion into two sections, answering both concluding questions; what was achieved in this thesis
and what might future developments be.

\section{Review}


This thesis introduced in Chapter (\ref{chapter: problem}) the newly proposed adaptive Deep
Orthogonal Decomposition based Deep Learning Reduced Order Model (DOD-DL-ROM)
based on the founding idea of \cite{franco2024} and a crossing with the POD-DL-ROM
given by \cite{fresca2022}. It further
concluded in its theoretical framework in Chapter (\ref{chapter: results}) 
of approximation theoretic statements, that
the Deep Orthogonal Decomposition based Reduced Order Model might pose a viable alternative to the 
Proper Orthogonal Decomposition based one. This is the case, as 
long as the aimed upon error tolerance $\varepsilon$ is significantly small, the FOM dimension $N_h$ is containable,
and enough regularity onto the optimal Deep Orthogonal Decomposition can be implied. The bottleneck
criterion for this is most significantly given by 
\begin{equation*}
    \frac{n}{s-1} > \frac{\log(N_h)}{\log(\varepsilon^{-1})}.
\end{equation*}
Here, $s$ denotes the implied regularity of the optimal non-linear embedding function between the coefficient
space and the latent dynamics.

We reached this conclusion through a rigorous application of results about arbitrary proximity to the
optimal projection
derived from the paper \cite{franco2024}. This happened in a context of approximation theory, lifted, in notation and
content, from \cite{gühring2019nn}. Furthermore, we could provide a distinct error decomposition upper and
lower bound for any such Deep Orthogonal Decomposition projection-based model. We went on to visualize this
in the last of three examples, whereas the other two served the purpose of showing one instance of potential
superiority of the method over the legacy Proper Orthogonal data-driven model, and another of its limitations.

In a different contribution, we have also provided a CoLoRA inspired architecture, that leverages a limit-fashion 
stationary version of the PPDE. We could not provide any meaningful numerical results for this. However, 
the model can be found in the repository \url{https://github.com/dawid-kotowski/doddlrom}.


\section{Research}


In the following, a few different hot spots of potential sources of interest will be highlighted, even though
the exact nature of their respective potential has not been investigated.
For this, we will list a few talking points regarding the proposed DOD-DL-ROM and any other
adaptations of the newly time-dependent DOD Definition (\ref{def: DOD}).
\begin{enumerate}
    \item[-] In Assumption (\ref{assumption additional lipschitz}), another imposition on the optimal DOD map
    has been made, and further even $W^{m, \infty}$-smoothness
    has been conditioned. There has been made no effort to show or discuss neither an example nor a criterion 
    to guarantee such regularity yet.
    \item[-] For the Upper Complexity Bounds, we have assumed, that $N_h = N_A$ to make the discussion regarding
    pre-reduction on the solution manifold less cumbersome, but this could be systematically reworked in a 
    broader class of problems, if the KnW behavior can be inherited through such a linear pre-reduction.
    \item[-] A simple comparison to a more direct approach like the older DL-ROM proposed by \cite{fresca2020} in
    terms of complexity.
    \item[-] In the same sense as the POD-DeepONet could be derived from the statement of the POD, there might
    be a possibility to undertake a generalization of the DOD projection to DeepONets in order to consider features slowing
    the KnW in a less strictly "parametric" way, leading to an adaptation of function space influence. 
    In such a way, the approach could be a candidate to overcome the \emph{curse of dimensionality}, if the
    features could be sampled in a dimension-independent way.
    \item[-] As it stands, the DOD uses $t \in [0, T]$ as a simple parameter and thus introduces time by
    injection as a factor of slow KnW decay. This might be not necessary in some cases. Hence, an alternative
    formulation might be of interest.
\end{enumerate}


\section*{Acknowledgement}

We hereby acknowledge and thank Björn Balkow and Felix Tümpner for review and help with notation.